\section{Integration of RAFFLE into Structure Generation}

\subsection{Motivation for RAFFLE Use: Enhanced Interface Diversity}

ARTEMIS-generated stacks, though consistent, impose alignment constraints. RAFFLE offers a mechanism to interpolate interfacial zones, increase disorder, or induce registry shifts.

\subsection{Current ARTEMIS-Only Pipeline: Limitations}

Aligned stacking methods limit structural exploration and may fail to capture intermediate configurations present in experimental systems.

\subsection{RAFFLE Capabilities: Refill Methods, Semi-Coherent Boundaries}

Using the \texttt{raffle\_generator} framework and reference energies from MACE, RAFFLE allows stochastic structural sampling within bounded interface zones. Visuals such as \texttt{upper\_a45\_b45\_c45.png} and \texttt{lower\_a45\_b45\_c45.png} represent generic orientation snapshots and are not RAFFLE-specific.

\subsection{Planned Integration Workflow}

\subsubsection{Combine ARTEMIS for stacking + RAFFLE for region refinement}

Initial slabs and orientations will be generated via ARTEMIS, while RAFFLE will introduce configurational variation within interfacial volumes.

\subsubsection{Target systems: strained, incoherent, or alloyed interfaces}

This approach is especially relevant to Si|Sn and Ge|Sn systems where lattice mismatch or registry discontinuity poses challenges for baseline methods.

\subsection{Validation and Comparison with ARTEMIS-Only Structures}

Relaxation comparisons will be made using both DFT and MLP pipelines to quantify RAFFLE's impact on interface energy diversity and ranking stability.

% TODO: Add schematic or figure comparing ARTEMIS-only and RAFFLE-refined structure sets